% ============================================================
% 1. Introduction
% ============================================================

\section{Introduction}
\label{sec:introduction}

Modern multi-agent LLM systems make an implicit promise of infinite
capacity: every sub-task spawns a new agent, each carrying its own
state---task descriptions, reasoning traces, embedding caches, tool
outputs---and the framework assumes that every agent's state will
remain available indefinitely.  In practice the resources run out.
LangGraph~\citep{langgraph2024}, AutoGen~\citep{wu2023autogen},
CrewAI~\citep{crewai2024}, and MetaGPT~\citep{hong2023metagpt} all
store agent states in unbounded Python data structures.  A recursive
research decomposition with branching factor~$k=3$ and depth~$d=7$
produces $\sum_{i=0}^{7}3^i = 3{,}280$ agents, each carrying 1--4\,KB
of LLM context; at depth~10 the count exceeds one million.  Without
explicit memory management, these systems crash.

The operating-systems community solved analogous problems decades ago
with virtual memory, page replacement, and cache hierarchies.  Yet the
agentic-AI community has largely ignored this body of work, relying on
ad-hoc solutions: MemGPT~\citep{packer2023memgpt} pages individual
agent context to disk but does not bound total memory;
Reflexion~\citep{shinn2023reflexion} compresses a single agent's
history but offers no multi-agent guarantees; Generative
Agents~\citep{park2023generative} maintains an unbounded memory stream.

We propose the \textbf{Bounded Infinity Cache (BIC)}, a memory-bounded
state-management primitive for recursive agent swarms.  BIC is not a
new memory-management paradigm; it is a specific instantiation that
combines (a)~the deterministic addressing of streaming algorithms,
(b)~the semantic compression of LLM-based summarization, and (c)~the
recursive agent state-management problem that operating-system
paging, streaming sketches, and RAG retrieval do not directly target
(\S\ref{sec:related}).  We use ``cache-bounded'' precisely: the cache
holding full agent states is bounded by $\Oh(M)$ for a fixed capacity
$M$, while a lightweight registry ($\sim$100 bytes per agent)
tracking tree metadata grows as $\Oh(N)$.  BIC draws on two core
mathematical tools:

\begin{enumerate}
  \item \textbf{Cantor pairing functions} map arbitrary tree-structured
        agent hierarchies to natural-number addresses, which are then
        reduced modulo $M$ to assign cache slots with structural
        locality (siblings land in nearby slots).
  \item \textbf{Hierarchical summarization} compresses evicted states
        into their parent's summary buffer via a lossy function
        $\summarize$, creating a chain of summaries up the tree that
        supports $\Oh(\log N)$ reconstruction of any evicted state.
\end{enumerate}

\noindent
We additionally explore \textbf{Hilbert space-filling curves} as an
optional locality-aware eviction optimization that compresses
high-dimensional agent-state vectors to one-dimensional indices,
aiming to cluster similar states for joint summarization.  However,
ablation experiments (\S\ref{sec:experiments}) show no measurable
benefit from Hilbert clustering in our current experimental settings.

\paragraph{Contributions.} We make three contributions:
\begin{enumerate}
  \item We introduce the BIC architecture---a specific instantiation
        in the lineage of operating-system memory management,
        streaming algorithms, and RAG/LLM-memory systems
        (\S\ref{sec:related}) that combines deterministic structural
        addressing with semantic-compression-based eviction for the
        recursive agent state problem.  We prove four theorems
        guaranteeing bounded cache consumption, zero fragmentation,
        information-preserving retrievability, and indefinite
        liveness (\S\ref{sec:method}).
  \item We design a hierarchical summarization scheme with
        Cantor-pairing-based addressing that preserves parent--child
        locality, and formally bound the information loss under
        repeated eviction at $(1-\epsilon)^d$ per $d$ summarization
        levels.  We additionally describe Hilbert-curve locality
        eviction as an optional optimization; ablation studies show
        it provides no measurable benefit in current experiments
        (\S\ref{sec:proofs}, \S\ref{sec:experiments}).
  \item We empirically demonstrate on synthetic and real-LLM
        (Gemini-2.0-flash) research-decomposition tasks that BIC
        achieves $100\%$ query success with $0.619$ reconstruction
        quality in bounded cache, outperforming six baselines
        (including LRU+summarization and a Tiered MemGPT-style
        baseline) by $32\times$ on reconstruction quality at scale,
        and $4.1\times$ on semantic quality with real LLM outputs
        at 20\% retention (\S\ref{sec:experiments}).
\end{enumerate}
